\documentclass[11pt]{article}
% Shared declarations; root documents always import a concrete semantic module.
\usepackage{business-admin}
\usepackage{amsmath,amssymb}
\usepackage{longtable,tabularx,array,booktabs}
\usepackage{enumitem}
\usepackage{ragged2e}
\newcolumntype{Y}{>{\RaggedRight\arraybackslash}X}
\newcommand{\SuiteMeta}[2]{\def\DocTitle{#1}\def\DocSubtitle{#2}\def\DocVersion{1.0}\def\DocStatus{Educational reference}\def\DocOwner{Small Business Management Suite}\def\DocSummary{Original teaching and practice resource; no real-business approval implied.}}
\newcommand{\SuiteOpen}{\makedoctitle\vspace{-0.7em}\noindent\rule{\linewidth}{0.45pt}\par\medskip}
\newcommand{\SourceNote}[2]{\par\medskip\noindent{\footnotesize\textbf{Reading anchor:} Longenecker, Petty, Palich and Hoy, \emph{Small Business Management: Launching \& Growing Entrepreneurial Ventures}, 20th ed., Ch.~#1, textbook p.~#2. This document is independently authored commentary and exercises, not a reproduction of the chapter.}}
\newcommand{\WritingLines}[1]{\par\smallskip\noindent\rule{\linewidth}{0.35pt}\par\vspace{0.35cm}\noindent\rule{\linewidth}{0.35pt}\par\vspace{0.35cm}\noindent\rule{\linewidth}{0.35pt}\par\medskip}
\newcommand{\FormField}[1]{\par\medskip\noindent\textbf{#1}\par\WritingLines{3}}
\newcommand{\SBMCheckBox}{\(\square\)}
\newcommand{\AssessmentRow}[2]{#1 & #2 & \(\square\) & \(\square\) & \(\square\) & \(\square\) & \(\square\)\\\hline}
\newcommand{\AssessmentHeader}{\noindent\begin{tabular}{|p{0.22\linewidth}|p{0.49\linewidth}|c|c|c|c|c|}\hline\textbf{Criterion} & \textbf{Evidence to inspect} & \textbf{0} & \textbf{1} & \textbf{2} & \textbf{3} & \textbf{4}\\\hline}
\newcommand{\AssessmentEnd}{\end{tabular}\par}
\setlength{\emergencystretch}{3em}

\SuiteMeta{Finance Lab: Pricing and Break-Even Analysis}{Original worked numbers | instructional only}
\begin{document}
\SuiteOpen
\section{Purpose and constraints}
Use the fictional figures provided. Show time units and formulas; make rounding explicit. Do not generalize prices, borrowing rates, margins, multiples or risk thresholds as representative of real-world market conditions.
\section{Inputs}
\begin{tabularx}{\linewidth}{@{}Y r l@{}}\toprule Input & Value & Unit\\\midrule
Proposed unit price & 75 & USD/unit\\
Variable cost & 27 & USD/unit\\
Period fixed costs & 9,600 & USD\\
Target volume & 250 & units\\
Alternative price & 69 & USD/unit\\
\bottomrule\end{tabularx}
\section{Calculation procedure}
\begin{enumerate}[leftmargin=*,itemsep=5pt]
\item State what quantity must be computed and over what period.
\item Select consistent source inputs and define each formula.
\item Compute the result, independently validate arithmetic, and show the cash or decision implication.
\item Challenge one assumption and report the changed decision.
\end{enumerate}

\section{Worked answer and checks}
\begin{tabularx}{\linewidth}{@{}p{0.30\linewidth}Y@{}}\toprule Measure & Calculation\\\midrule
Unit contribution & 75 - 27 = 48\\
Break-even volume & 9600 / 48 = 200 units\\
Profit at 250 units & 250 x 48 - 9600 = 2400\\
Break-even after price reduction & ceil(9600 / (69 - 27)) = 229 units\\
\bottomrule\end{tabularx}
\section{Extension exercise}
Determine the increase in demand required for the lower price to outperform the original 250-unit profit; distinguish markup from contribution margin.\par\textbf{Required answer:} state calculations, evidence gaps, alternative scenarios and next data-collection steps.
\section{Review checklist}
\begin{itemize}[leftmargin=*,itemsep=3pt]
\item Inputs, rate conventions and units are explicit.
\item Recalculations match all published figures within stated rounding.
\item Profit, cash and valuation are not confused.
\item The recommendation includes adverse-scenario and decision thresholds.
\end{itemize}

\end{document}
